\chapter{Complexity — Hierarchy}
%%% generated by the blueprint pipeline from TCSlib.Complexity.CircuitComplexity.Hierarchy
\section{Overview}

Declarations of the Lean module \texttt{TCSlib.Complexity.CircuitComplexity.Hierarchy}.

\section{Declarations}

\begin{theorem}[Value of an empty gate]
\lean{BoolCircuit.eval_node_nil}
\label{BoolCircuit.eval_node_nil}
\leanok
\uses{BoolCircuit.TreeCircuit, BoolCircuit.TreeCircuit.eval}
\difficulty{2}
For any number $n$ of Boolean variables, any assignment to them, and either choice of
gate, the value of a gate with no children is true for an AND gate and false for an OR
gate; equivalently, its value is the Boolean value specifying the gate choice.
\end{theorem}

\begin{definition}[Widening a Boolean circuit]
\lean{BoolCircuit.widenCircuit}
\label{BoolCircuit.widenCircuit}
\leanok
\uses{BoolCircuit.TreeCircuit}
Let $m\le n$. A circuit on $m$ variables is widened to one on $n$ variables by
preserving each gate and its structure, and replacing every literal's variable index by
the same index viewed among the $n$ variables, with its polarity unchanged.
\end{definition}

\begin{theorem}[Evaluation under a wider input space]
\lean{BoolCircuit.widenCircuit_eval}
\label{BoolCircuit.widenCircuit_eval}
\leanok
\uses{BoolCircuit.Lit.eval, BoolCircuit.TreeCircuit, BoolCircuit.TreeCircuit.eval, BoolCircuit.TreeCircuit.ind, BoolCircuit.widenCircuit}
\difficulty{4}
Let $m\text{ and }n$ be nonnegative integers with $m\le n$, let $c$ be a Boolean circuit
tree on $m$ variables, and let $x$ assign Boolean values to the $n$ variables. The
circuit obtained by widening $c$ to $n$ variables evaluates under $x$ to the value of
$c$ under the assignment that sends each index $i<m$ to the value of $x$ at the
corresponding index in the larger input space.
\end{theorem}

\begin{theorem}[Size preserved by widening]
\lean{BoolCircuit.widenCircuit_size}
\label{BoolCircuit.widenCircuit_size}
\leanok
\uses{BoolCircuit.TreeCircuit, BoolCircuit.TreeCircuit.ind, BoolCircuit.TreeCircuit.size, BoolCircuit.widenCircuit}
\difficulty{4}
Let $m\text{ and }n$ be nonnegative integers with $m\le n$, and let $c$ be a Boolean
circuit tree on $m$ variables. Widening the circuit to $n$ variables preserves its total
number of nodes.
\end{theorem}

\begin{definition}[Extending a Boolean assignment]
\lean{BoolCircuit.extendBy}
\label{BoolCircuit.extendBy}
\leanok
Let $m,n$ be natural numbers, let $\mathrm{pad}$ assign a Boolean value to each of the
$n$ indices, and let $y$ assign a Boolean value to each of the $m$ indices. Define an
assignment on the $n$ indices by using $y(i)$ when $i<m$, and using $\mathrm{pad}(i)$
otherwise.
\end{definition}

\begin{definition}[Restriction of a Boolean circuit]
\lean{BoolCircuit.restrictCircuit}
\label{BoolCircuit.restrictCircuit}
\leanok
\uses{BoolCircuit.Lit.eval, BoolCircuit.TreeCircuit}
For natural numbers $m,n$, an assignment $p : \mathrm{Fin}\,n \to \mathrm{Bool}$, and a
Boolean circuit $C$ on $n$ variables, its restriction to $m$ variables is the circuit
defined recursively as follows: a literal with index less than $m$ is retained with that
index viewed among the $m$ variables, while a literal with index at least $m$ is
replaced by a node with no children whose tag is the value of the literal under $p$.
Each internal node keeps its tag, and the restriction is applied to each child.
\end{definition}

\begin{theorem}[Evaluation after circuit restriction]
\lean{BoolCircuit.restrictCircuit_eval}
\label{BoolCircuit.restrictCircuit_eval}
\leanok
\uses{BoolCircuit.Lit.eval, BoolCircuit.TreeCircuit, BoolCircuit.TreeCircuit.eval, BoolCircuit.TreeCircuit.ind, BoolCircuit.eval_node_nil, BoolCircuit.extendBy, BoolCircuit.restrictCircuit}
\difficulty{5}
Let $m,n\in\mathbb{N}$, let $p$ be an assignment of Boolean values to the $n$ variables,
and let $C$ be a Boolean circuit on those variables. For every assignment $y$ to the
first $m$ variables, restricting $C$ to those variables and replacing each literal whose
index is at least $m$ by its value under $p$ gives a circuit whose value at $y$ equals
the value of $C$ on the assignment that uses $y_i$ for indices $i<m$ and $p_i$ for
indices $i\ge m$.
\end{theorem}

\begin{theorem}[Size preservation under circuit restriction]
\lean{BoolCircuit.restrictCircuit_size}
\label{BoolCircuit.restrictCircuit_size}
\leanok
\uses{BoolCircuit.TreeCircuit, BoolCircuit.TreeCircuit.ind, BoolCircuit.TreeCircuit.size, BoolCircuit.restrictCircuit}
\difficulty{4}
Let $m,n$ be natural numbers, let $p : \\mathrm{Fin}\,n \\to \\bbf_2$ be a Boolean
assignment, and let $C$ be a Boolean circuit on $n$ variables. Restrict each literal
whose variable index is less than $m$ to the corresponding variable among the first $m$
variables; replace each other literal by a constant gate with value given by that
literal evaluated at $p$, and retain the circuit's gate structure. The resulting circuit
has the same size as $C$.
\end{theorem}

\begin{definition}[Restriction to the first variables]
\lean{BoolCircuit.onFirst}
\label{BoolCircuit.onFirst}
\leanok
Given natural numbers $m \le n$ and a Boolean function $f$ on $m$ Boolean inputs, define
its extension to $n$ inputs by applying $f$ to the first $m$ coordinates of each input.
\end{definition}

\begin{theorem}[Restriction under dependence on initial variables]
\lean{BoolCircuit.restrictCircuit_eval_of_onFirst}
\label{BoolCircuit.restrictCircuit_eval_of_onFirst}
\leanok
\uses{BoolCircuit.TreeCircuit, BoolCircuit.TreeCircuit.eval, BoolCircuit.extendBy, BoolCircuit.onFirst, BoolCircuit.restrictCircuit, BoolCircuit.restrictCircuit_eval}
\difficulty{3}
Let $m\\le n$, let $p$ be an assignment to $n$ variables, and let $C$ be a circuit on
$n$ variables whose evaluation on every assignment $x$ equals $f$ applied to the
restriction of $x$ to its first $m$ variables. Then, for every assignment $y$ to those
$m$ variables, the circuit obtained by restricting $C$ to the first $m$ variables and
using $p$ for the remaining variables evaluates to $f(y)$.
\end{theorem}

\begin{definition}[Languages recognized within a tree-size bound]
\lean{Language.InTreeSize}
\label{Language.InTreeSize}
\leanok
\uses{BoolCircuit.TreeCircuit, BoolCircuit.TreeCircuit.eval, BoolCircuit.TreeCircuit.size}
For a function $T:\mathbb{N}\to\mathbb{N}$, a language $L$ of Boolean words has tree
size bounded by $T$ if there is a family of Boolean tree circuits
$(C_n)_{n\in\mathbb{N}}$, with $C_n$ acting on $n$ variables, such that
$\mathrm{size}(C_n)\le T(n)$ for every $n$, and every Boolean word $w$ belongs to $L$
exactly when $C_{|w|}$ evaluates to true on the coordinates of $w$.
\end{definition}

\begin{theorem}[Monotonicity of tree circuit size]
\lean{Language.InTreeSize.mono}
\label{Language.InTreeSize.mono}
\leanok
\uses{Language.InTreeSize}
\difficulty{3}
Let $T,T':\mathbb{N}\to\mathbb{N}$ and let $L$ be a language of Boolean words. If $L$
has Boolean circuit trees of size at most $T(n)$ that recognize it at each input length,
and $T(n)\le T'(n)$ for every $n$, then $L$ also has such trees of size at most $T'(n)$
at every input length.
\end{theorem}

\begin{theorem}[The empty language has bounded tree size]
\lean{Language.zero_inTreeSize}
\label{Language.zero_inTreeSize}
\leanok
\uses{BoolCircuit.TreeCircuit.size, BoolCircuit.eval_node_nil, Language.InTreeSize}
\difficulty{3}
Let $T$ be a function from the natural numbers to the natural numbers such that $1\le
T(n)$ for every $n$. The empty language has a Boolean tree circuit for each input length
$n$, of size at most $T(n)$, whose acceptance condition agrees with membership in the
empty language on every Boolean string.
\end{theorem}

\begin{definition}[Languages of bounded tree size]
\lean{BoolCircuit.TreeSize}
\label{BoolCircuit.TreeSize}
\leanok
\uses{Language.InTreeSize}
Given a function $T:\mathbb{N}\to\mathbb{N}$, the set of languages of Boolean words
having tree size bounded by $T$ consists of exactly those languages recognized by a
family of Boolean tree circuits whose size at each input length $n$ is at most $T(n)$.
\end{definition}

\begin{theorem}[Characterization of tree-size languages]
\lean{BoolCircuit.mem_treeSize_iff}
\label{BoolCircuit.mem_treeSize_iff}
\leanok
\uses{BoolCircuit.TreeSize, Language.InTreeSize}
\difficulty{1}
For a size bound $T : \mathbb{N} \to \mathbb{N}$ and a language $L$ over Boolean words,
$L$ belongs to the class of languages of tree size at most $T$ if and only if there is a
family of Boolean circuit trees, one on $n$ variables for each $n$, such that each tree
has size at most $T(n)$ and, for every Boolean word $w$, membership of $w$ in $L$ is
equivalent to the tree on $|w|$ variables evaluating to true on the assignment given by
the entries of $w$.
\end{theorem}

\begin{theorem}[Equality of circuit evaluations]
\lean{BoolCircuit.eval_eq_of_iff}
\label{BoolCircuit.eval_eq_of_iff}
\leanok
\uses{BoolCircuit.TreeCircuit, BoolCircuit.TreeCircuit.eval}
\difficulty{4}
Let $C_n$ and $D_n$ be tree circuits on $n$ inputs for each natural number $n$. Suppose
that for every Boolean list $w$, the circuits $C_{|w|}$ and $D_{|w|}$ have the same
truth value under the assignment given by the entries of $w$. Then for every natural
number $n$ and every Boolean assignment to $n$ inputs, $C_n$ and $D_n$ have the same
truth value.
\end{theorem}

\begin{definition}[Padding a circuit family]
\lean{BoolCircuit.padFamily}
\label{BoolCircuit.padFamily}
\leanok
\uses{BoolCircuit.TreeCircuit, BoolCircuit.universalCircuit, BoolCircuit.widenCircuit}
Let $\ell:\mathbb{N}\to\mathbb{N}$ satisfy $\ell(n)\le n$ for every $n$, and let $F_n$
be a Boolean function on $\ell(n)$ inputs for each $n$. For each $n$, define an
$n$-variable Boolean tree circuit by extending the universal circuit for $F_n$ to $n$
variables.
\end{definition}

\begin{theorem}[Evaluation of a padded circuit family]
\lean{BoolCircuit.padFamily_eval}
\label{BoolCircuit.padFamily_eval}
\leanok
\uses{BoolCircuit.TreeCircuit.eval, BoolCircuit.onFirst, BoolCircuit.padFamily, BoolCircuit.universalCircuit_eval, BoolCircuit.widenCircuit_eval}
\difficulty{3}
Let $\ell:\mathbb{N}\to\mathbb{N}$ satisfy $\ell(n)\le n$ for every $n$, and let $F_n$
be a Boolean function on assignments to the first $\ell(n)$ inputs. For every $n$ and
every assignment $x$ to $n$ inputs, the padded circuit for $F_n$ evaluates to $F_n$ on
the restriction of $x$ to its first $\ell(n)$ inputs.
\end{theorem}

\begin{definition}[Language recognized by padded circuits]
\lean{BoolCircuit.padLanguage}
\label{BoolCircuit.padLanguage}
\leanok
\uses{BoolCircuit.TreeCircuit.eval, BoolCircuit.padFamily}
Let $\ell:\mathbb{N}\to\mathbb{N}$ satisfy $\ell(n)\le n$ for every $n$, and let $F_n$
be a Boolean function on $\ell(n)$ inputs for each $n$. The language of Boolean words
$w$ for which the padded circuit at length $|w|$ evaluates to true on $w$ is the
language recognized by this family.
\end{definition}

\begin{theorem}[Tree-size representation of a padded language]
\lean{BoolCircuit.padLanguage_inTreeSize}
\label{BoolCircuit.padLanguage_inTreeSize}
\leanok
\uses{BoolCircuit.padFamily, BoolCircuit.padLanguage, BoolCircuit.universalCircuit_size_le, BoolCircuit.widenCircuit_size, Language.InTreeSize}
\difficulty{3}
Let $\ell,T':\mathbb{N}\to\mathbb{N}$ satisfy $\ell(n)\le n$ for every $n$, and let
$F_n$ be a Boolean function on assignments to $\ell(n)$ inputs. Suppose that for every
$n$, $2^{\ell(n)}(\ell(n)+1)+1\le T'(n)$. Then the language of Boolean lists accepted by
the padded circuits has a family of tree circuits, one on each input length, whose size
at length $n$ is at most $T'(n)$ and which accepts exactly the lists in that language.
\end{theorem}

\begin{theorem}[Failure of bounded tree-size representation]
\lean{BoolCircuit.padLanguage_not_inTreeSize}
\label{BoolCircuit.padLanguage_not_inTreeSize}
\leanok
\uses{BoolCircuit.TreeCircuit, BoolCircuit.TreeCircuit.eval, BoolCircuit.TreeCircuit.size, BoolCircuit.eval_eq_of_iff, BoolCircuit.onFirst, BoolCircuit.padFamily_eval, BoolCircuit.padLanguage, BoolCircuit.restrictCircuit, BoolCircuit.restrictCircuit_eval_of_onFirst, BoolCircuit.restrictCircuit_size, Language.InTreeSize}
\difficulty{4}
Let $\ell,T:\mathbb{N}\to\mathbb{N}$, let $n_0$ be a natural number, and suppose
$\ell(n)\le n$ for every $n$. Let $F_n$ be a Boolean function on assignments to
$\ell(n)$ inputs. Assume that every tree circuit on $\ell(n_0)$ inputs of size at most
$T(n_0)$ disagrees with $F_{n_0}$ on some assignment. Then the language accepted by the
padded circuits has no family of tree circuits whose size at length $n$ is at most
$T(n)$ and which recognizes that language exactly.
\end{theorem}

\begin{theorem}[Strict Inclusion of Tree Size Classes]
\lean{BoolCircuit.treeSize_ssubset}
\label{BoolCircuit.treeSize_ssubset}
\leanok
\uses{BoolCircuit.TreeCircuit, BoolCircuit.TreeCircuit.eval, BoolCircuit.TreeCircuit.size, BoolCircuit.TreeSize, BoolCircuit.exists_not_eval_of_lt, BoolCircuit.padLanguage_inTreeSize, BoolCircuit.padLanguage_not_inTreeSize, Language.InTreeSize.mono}
\difficulty{5}
Let $T,T',\ell:\mathbb N\to\mathbb N$ and let $n_0\in\mathbb N$. Suppose that
$\ell(n)\le n$ and $T(n)\le T'(n)$ for every $n$, that $2^{\ell(n)}(\ell(n)+1)+1\le
T'(n)$ for every $n$, and that $(\ell(n_0)+4)T(n_0)<2^{\ell(n_0)}$. Then the class of
Boolean languages recognizable by circuit-tree families of size at most $T(n)$ on inputs
of length $n$ is a proper subset of the corresponding class with size bound $T'(n)$.
\end{theorem}

\begin{theorem}[Strict Inclusion Under a Size Gap]
\lean{BoolCircuit.treeSize_ssubset_of_lt}
\label{BoolCircuit.treeSize_ssubset_of_lt}
\leanok
\uses{BoolCircuit.TreeSize, BoolCircuit.treeSize_ssubset}
\difficulty{2}
Let $T:\mathbb N\to\mathbb N$ and $n_0\in\mathbb N$ satisfy $(n_0+4)T(n_0)<2^{n_0}$.
Then the class of Boolean languages recognizable by circuit-tree families of size at
most $T(n)$ on inputs of length $n$ is a proper subset of the class recognizable with
size at most $\max\bigl(T(n),2^{\min(n,n_0)}(\min(n,n_0)+1)+1\bigr)$ on inputs of length
$n$.
\end{theorem}

\begin{theorem}[Strict Inclusion of Tree Size Classes]
\lean{BoolCircuit.treeSize_one_ssubset}
\label{BoolCircuit.treeSize_one_ssubset}
\leanok
\uses{BoolCircuit.TreeSize, BoolCircuit.treeSize_ssubset_of_lt}
\difficulty{2}
The class of Boolean languages that can be recognized by a family of Boolean circuit
trees whose size is at most $1$ for every input length is a proper subset of the class
recognizable by such a family with size at most
$\max\bigl(1,2^{\min(n,3)}(\min(n,3)+1)+1\bigr)$ on inputs of length $n$.
\end{theorem}

\begin{theorem}[Empty language of tree size one]
\lean{BoolCircuit.zero_mem_treeSize_one}
\label{BoolCircuit.zero_mem_treeSize_one}
\leanok
\uses{BoolCircuit.TreeSize, Language.zero_inTreeSize}
\difficulty{1}
The empty language over the Boolean alphabet has tree size at most one: there is a
family of Boolean tree circuits, one for each input length, each of size at most $1$,
whose accepted language is empty.
\end{theorem}

\begin{theorem}[Tree to fan-in-two circuit size bound]
\lean{Language.InTreeSize.inSIZE}
\label{Language.InTreeSize.inSIZE}
\leanok
\uses{BoolCircuit.DAGCircuit.size, BoolCircuit.DAGCircuitFamily.mem_language_iff, BoolCircuit.TreeCircuit.size, BoolCircuit.TreeCircuit.toBinary, BoolCircuit.TreeCircuit.toDAG, BoolCircuit.TreeCircuit.toDAG_eval, BoolCircuit.TreeCircuit.toDAG_isFaninTwo, BoolCircuit.TreeCircuit.toDAG_size_le, BoolCircuit.toBinary_eval, BoolCircuit.toBinary_maxFanin_le, BoolCircuit.toBinary_size_le, Language.InSIZE, Language.InTreeSize}
\difficulty{4}
Let $T:\mathbb{N}\to\mathbb{N}$ and let $L$ be a language of Boolean words. If for every
input length $n$ there is a Boolean circuit tree of size at most $T(n)$ that accepts
exactly the words of length $n$ in $L$, then there is a family of well-formed acyclic
Boolean circuits with fan-in at most two that recognizes $L$ and has size at most
$n+3T(n)$ on inputs of length $n$.
\end{theorem}
